% GENERATED FILE. Edit canonical lessons, then run npm run generate:books.
% canonical-source-hash: fnv1a64:b508869b23439594
% canonical-lessons: GE-C33-ja-nein-frage, GE-C33-imperativ, GE-C33-wer, GE-C33-wann, GE-C33-warum, GE-R33-frag-mich

\chapter{When the Verb Steps to the Front}
\label{ch:ge-frage}
\hlchaptermodality{46}

\begin{chapteropening}
\textbf{By the end of this chapter:} \emph{I can ask a question that expects yes or no, give a one-word command, and ask who, when and why.}

Every statement in the book becomes a question and every verb becomes a command. Three of the six W-words were already owned, so the chapter buys the whole interrogative system for three new words plus two sentence shapes -- and the fourth, fifth and sixth W arrive predicted rather than memorised, because Grimm's law lines German w- up against English wh- one word at a time.
\end{chapteropening}

\section[Wohnst du in Berlin?]{Wohnst du in Berlin? — the question that has no question word}

\label{lesson:GE-C33-ja-nein-frage}



\begin{quote}
You have owned \textbf{ja}, \textbf{nein} and \textbf{doch} since the yes-and-no chapter, and you have never been shown the question they answer. Say the three first, then read on.

\begin{itemize}
\raggedright
  \item \emph{Recall:} two chapters back, the two negatives and what each cancels — \emph{kein} for a noun, \emph{nicht} for the rest
\end{itemize}
\end{quote}

\begin{grammarlens}[title={Grammar Lens: the verb moves to the front}]
\begin{quote}
To ask something answerable with \emph{ja} or \emph{nein}, take the statement and put the \textbf{verb first}.
\end{quote}

Nothing else changes. No extra word, no rising list of rules:

\begin{quote}
\textbf{Du wohnst in Berlin.} $\to$ \textbf{Wohnst du in Berlin?} · \textbf{Er hat einen Bruder.} $\to$ \textbf{Hat er einen Bruder?}
\end{quote}

The verb has stood in slot two for the whole book. Move it one place left and the sentence becomes a question.
\end{grammarlens}

\begin{culture}
English used to do this and stopped. \emph{\textquotedblleft{}Speak you German\textquotedblright{}} is ordinary Elizabethan word order. Somewhere after Shakespeare, English stopped moving its verbs and built a machine instead — the \textbf{do} in \emph{do you speak German}, which means nothing at all. German never built it, because it never lost the habit. The German question is the older one; the English \emph{do} is the innovation.
\end{culture}

\subsection*{Your turn}
Say the statement, then flip it.

Say these aloud:

\begin{itemize}
\raggedright
  \item \textquotedblleft{}Du wohnst in Berlin.\textquotedblright{} $\to$ \textquotedblleft{}Wohnst du in Berlin?\textquotedblright{}
  \item \textquotedblleft{}Er hat einen Bruder.\textquotedblright{} $\to$ \textquotedblleft{}Hat er einen Bruder?\textquotedblright{}
  \item the three drinks, asked for — \textquotedblleft{}Hast du Kaffee, Tee, Milch?\textquotedblright{}
  \item answer with last chapter's negative — \textquotedblleft{}Nein, ich habe keinen Kaffee.\textquotedblright{}
  \item the contradiction — a negative question, answered \textquotedblleft{}Doch!\textquotedblright{}
\end{itemize}

\emph{Twice through:} \textquotedblleft{}Wohnst du in Berlin? — Ja. / Nein.\textquotedblright{}

\begin{tcolorbox}[breakable,colback=teal!4,colframe=teal!35!black,title={Before you move on}]
What turns a German statement into a yes-or-no question? (\textbf{The verb moves to first place.}) What did English build when it stopped? (The empty \textbf{do}.) Which answer-word contradicts a negative question? (\emph{\textbf{Doch}}.)
\end{tcolorbox}
\section[Komm!]{Komm! — the sentence with no subject at all}

\label{lesson:GE-C33-imperativ}



\begin{quote}
Take \textbf{kommen} and strip the \emph{-en} off it. What is left is the stem, \textbf{komm-}, and it is about to be a whole sentence on its own.
\end{quote}

\begin{grammarlens}[title={Grammar Lens: the stem, alone}]
\begin{quote}
To tell one person you are on \emph{du} terms with to do something, say the \textbf{stem and nothing else}.
\end{quote}

\begin{quote}
\textbf{Komm!} — \textquotedblleft{}Come!\textquotedblright{} · \textbf{Geh!} — \textquotedblleft{}Go!\textquotedblright{} · \textbf{Sag!} — \textquotedblleft{}Say!\textquotedblright{}
\end{quote}

Two things are missing, and both matter. There is \textbf{no ending} — this is the one shape of a German verb that wears none. And there is \textbf{no pronoun}, in a language that has spent forty chapters insisting the subject slot be filled.

A command is the single place German lets that slot stay empty, because the person addressed is standing right there.
\end{grammarlens}

\begin{cousinweb}
The bare stem is the oldest shape a verb has anywhere in this family. English kept it too — \emph{come}, \emph{go}, \emph{say}, \emph{open}, \emph{listen} — which is why an English command and a German one line up word for word and syllable for syllable:

\begin{quote}
\textbf{Hör!} — \emph{Hear!} · \textbf{Sieh!} — \emph{See!} · \textbf{Lies!} — \emph{Read!} · \textbf{Schreib!} — \emph{Write!}
\end{quote}

Two of those hide something you already know: the vowel of \emph{sehen} and \emph{lesen} breaks in the \emph{du} form, and the command takes the broken vowel with it. \textbf{Sieh} and \textbf{lies}, not \emph{seh} and \emph{les}.
\end{cousinweb}

\subsection*{Your turn}
Find the stem before you say the command.

Say these aloud:

\begin{itemize}
\raggedright
  \item kommen $\to$ \textquotedblleft{}Komm!\textquotedblright{}
  \item gehen $\to$ \textquotedblleft{}Geh!\textquotedblright{}
  \item sagen $\to$ \textquotedblleft{}Sag!\textquotedblright{}
  \item the pair that opens and closes — \textquotedblleft{}Öffne!\textquotedblright{} and \textquotedblleft{}Schließ!\textquotedblright{}
  \item the one that hears — \textquotedblleft{}Hör!\textquotedblright{}
  \item the two with a broken vowel — \textquotedblleft{}Lies!\textquotedblright{} and \textquotedblleft{}Sieh!\textquotedblright{}
\end{itemize}

\emph{Twice through:} \textquotedblleft{}Komm! Geh! Hör!\textquotedblright{}

\begin{tcolorbox}[breakable,colback=teal!4,colframe=teal!35!black,title={Before you move on}]
What shape is a German command? (\textbf{The bare stem} — no ending.) What does it leave out that German otherwise demands? (\textbf{The pronoun.}) Both this and a yes-or-no question begin with the verb — so what tells them apart? (The question \textbf{has a subject after the verb}; the command has none.) Next: the six question words, and the letter they all begin with.
\end{tcolorbox}
\section[wer]{wer — \textquotedblleft{}who,\textquotedblright{} and the letter the whole family shares}

\label{lesson:GE-C33-wer}



\begin{quote}
You already own three question words — \textbf{was}, \textbf{wo}, \textbf{wie}. Say them, and notice what all three start with.
\end{quote}

\subsection*{You'll want to know: wer}
\begin{quote}
\textbf{wer} — \textquotedblleft{}who.\textquotedblright{}
\end{quote}

Said \emph{vair}: the \textbf{w} is a \emph{v}, as it has been since \emph{wie} and \emph{wir}, and the \textbf{e} is long.

\begin{quote}
\textbf{Wer ist das?} — \textquotedblleft{}Who is that?\textquotedblright{}
\end{quote}

Take care with one trap. German \textbf{wer} means \emph{who}; German \textbf{wo} means \emph{where}, and an English eye reads it as \emph{who}. The two swap appearances across the languages, and only practice fixes it.

\begin{cousinweb}
\textbf{wer} is Latin \textbf{quis} and English \textbf{who}, from Proto-Indo-European \emph{\textbf{*k\textsuperscript{w}is}}. Grimm's law turned \textbf{k\textsuperscript{w}} into Germanic \textbf{hw}; English wrote \emph{hw} and flipped it to \textbf{wh}; German wore the \emph{h} off. That one step lines the two languages up letter for letter:

\noindent
\begin{tabularx}{\linewidth}{@{}>{\raggedright\arraybackslash}X>{\raggedright\arraybackslash}X@{}}
\toprule
\textbf{German} & \textbf{English} \\
\midrule
\textbf{w}er & \textbf{wh}o \\
\textbf{w}as & \textbf{wh}at \\
\textbf{w}o & \textbf{wh}ere \\
\textbf{w}ann & \textbf{wh}en \\
\textbf{w}arum & \textbf{wh}y \\
\bottomrule
\end{tabularx}

Five of those six you can now predict rather than memorise, and the sixth breaks the pattern in both languages at once: German \textbf{wie} against English \textbf{how}, where English lost the \emph{w} instead of the \emph{h}.
\end{cousinweb}

\subsection*{Your turn}
Say these aloud:

\begin{itemize}
\raggedright
  \item \textquotedblleft{}wer\textquotedblright{} — \emph{vair}, and \textquotedblleft{}Wer ist das?\textquotedblright{}
  \item the trap, both ways — \textquotedblleft{}wer\textquotedblright{} is \emph{who}, \textquotedblleft{}wo\textquotedblright{} is \emph{where}
  \item about the people you know — \textquotedblleft{}Wer ist der Freund? Wer ist die Freundin?\textquotedblright{}
  \item \textquotedblleft{}Wer ist die Mutter? Wer ist der Vater?\textquotedblright{}
  \item the command from the last lesson, then the question — \textquotedblleft{}Sag!\textquotedblright{} then \textquotedblleft{}Wer ist das?\textquotedblright{}
\end{itemize}

\emph{Twice through:} \textquotedblleft{}Wer ist das?\textquotedblright{}

\begin{tcolorbox}[breakable,colback=teal!4,colframe=teal!35!black,title={Before you move on}]
What does \emph{wer} mean, and what is the trap? (\emph{\textbf{Who}} — and \emph{wo}, which looks like it, means \emph{where}.) Which Latin word is its cousin? (\emph{\textbf{Quis}}.) What sound law turns English \emph{wh-} into German \emph{w-}? (\textbf{Grimm's law}, which made \emph{k\textsuperscript{w}} into \emph{hw}; English kept the \emph{h}, German dropped it.) Next: the fifth W.
\end{tcolorbox}
\section[wann]{wann — \textquotedblleft{}when,\textquotedblright{} and the question the calendar was waiting for}

\label{lesson:GE-C33-wann}



\begin{quote}
Predict this one before you read it. English \emph{when} begins \emph{wh-}. What must the German be?
\end{quote}

\subsection*{You'll want to know: wann}
\begin{quote}
\textbf{wann} — \textquotedblleft{}when.\textquotedblright{}
\end{quote}

Said \emph{vann}, with a short flat \emph{a} and the \emph{v}-sound on the front. The double \textbf{nn} keeps the vowel short, which is the same spelling habit you met in \emph{Mann} and \emph{kann}.

\begin{quote}
\textbf{Wann ist das?} — \textquotedblleft{}When is that?\textquotedblright{}
\end{quote}

The clock chapter gave you \textbf{Es ist ein Uhr} and no way to ask for it. Now you have one:

\begin{quote}
\textbf{Wann?} — \textbf{Um ein Uhr.}
\end{quote}

\begin{cousinweb}
\textbf{wann} and English \textbf{when} are one word, on the same \emph{\textbf{*k\textsuperscript{w}-}} root that gave \emph{wer} and \emph{who}. So the second row of the table is filled in by the rule rather than by memory, and the vocabulary this book spent five chapters on is suddenly askable:

\begin{quote}
\textbf{Wann? — Am Montag.} · \textbf{Wann? — Im Januar.} · \textbf{Wann? — Im Dezember.}
\end{quote}

Seven weekdays, twelve months and four seasons have been sitting in your head with no question attached to them. One word attaches all twenty-three.
\end{cousinweb}

\subsection*{Your turn}
Say these aloud:

\begin{itemize}
\raggedright
  \item \textquotedblleft{}wann\textquotedblright{} — \emph{vann}, short \emph{a}
  \item \textquotedblleft{}Wann ist das?\textquotedblright{}
  \item two weekdays as answers — \textquotedblleft{}Am Montag\textquotedblright{}, \textquotedblleft{}Am Samstag\textquotedblright{}
  \item the day the sun owns — \textquotedblleft{}Am Sonntag\textquotedblright{}
  \item two months — \textquotedblleft{}Im Januar\textquotedblright{}, \textquotedblleft{}Im Dezember\textquotedblright{}
  \item the four seasons — \textquotedblleft{}im Frühling, im Sommer, im Herbst, im Winter\textquotedblright{}
\end{itemize}

\emph{Twice through:} \textquotedblleft{}Wann? — Am Montag.\textquotedblright{}

\begin{tcolorbox}[breakable,colback=teal!4,colframe=teal!35!black,title={Before you move on}]
Say \textquotedblleft{}when\textquotedblright{} in German and name the English word it is. (\emph{\textbf{Wann}}; \emph{when}.) Which root does it share with \emph{wer}? (\emph{\textbf{*k\textsuperscript{w}-}}.) Answer \emph{wann} with a weekday and with a month. (\textbf{Am Montag. Im Januar.}) Which season is English \emph{harvest}? (\emph{\textbf{Herbst}}.) Next: the sixth W, and the two words hiding inside it.
\end{tcolorbox}
\section[warum]{warum — \textquotedblleft{}why,\textquotedblright{} and the two words welded inside it}

\label{lesson:GE-C33-warum}



\begin{quote}
Five W-words are yours: \textbf{wer, was, wo, wann, wie}. Say all five, then take the sixth.

\begin{itemize}
\raggedright
  \item \emph{Recall:} two chapters back, the negative for a \emph{der}-word you do not have — \emph{keinen}
\end{itemize}
\end{quote}

\subsection*{You'll want to know: warum}
\begin{quote}
\textbf{warum} — \textquotedblleft{}why.\textquotedblright{}
\end{quote}

Said \emph{va-ROOM}, and the stress is on the \textbf{second} syllable — which is unusual enough to be worth marking, because almost every German word you have met so far leans on its first.

\begin{quote}
\textbf{Warum?} — \textquotedblleft{}Why?\textquotedblright{}
\end{quote}

\begin{quote}
\textbf{Warum lernst du Deutsch?} — \textquotedblleft{}Why are you learning German?\textquotedblright{}
\end{quote}

\begin{cousinweb}
\textbf{warum} is two words that stopped being two. The first half, \textbf{war-}, is an old shape of \textbf{was}; the second, \textbf{-um}, is the ordinary German \textbf{um}, \textquotedblleft{}about\textquotedblright{}. \emph{About what.}

English built its own \emph{why}-word the identical way and it is still readable: \textbf{wherefore} is \emph{where} plus \emph{for}, \textquotedblleft{}for what\textquotedblright{}. Shakespeare's \emph{\textquotedblleft{}wherefore art thou Romeo\textquotedblright{}} is not asking where Romeo is. It is asking \textbf{why}.

So both languages fused a question word to a small preposition, and both ended up with a \emph{why} that once meant \emph{about what} or \emph{for what}. German's fusion went further and hid the seam; English's is still visible, which is why \emph{wherefore} looks quaint and \emph{warum} looks like a single word.
\end{cousinweb}

\subsection*{Your turn}
Say these aloud:

\begin{itemize}
\raggedright
  \item \textquotedblleft{}warum\textquotedblright{} — \emph{va-ROOM}, stress on the second half
  \item the two pieces, then the word — \textquotedblleft{}war … um … warum\textquotedblright{}
  \item \textquotedblleft{}Warum lernst du Deutsch?\textquotedblright{}
  \item the whole set, in one breath — \textquotedblleft{}wer, was, wo, wann, warum, wie\textquotedblright{}
  \item ask about the drink — \textquotedblleft{}Warum Tee?\textquotedblright{}
\end{itemize}

\emph{Twice through:} \textquotedblleft{}Warum lernst du Deutsch?\textquotedblright{}

\begin{tcolorbox}[breakable,colback=teal!4,colframe=teal!35!black,title={Before you move on}]
Which syllable of \emph{warum} is stressed? (\textbf{The second} — \emph{va-ROOM}.) What two words are welded inside it? (An old \textbf{was} and \textbf{um}, \textquotedblleft{}about\textquotedblright{}.) Which English word was built the same way? (\emph{\textbf{Wherefore}} — \emph{where} plus \emph{for}.) Say all six W-words. Next: the chapter, run cold.
\end{tcolorbox}
\section[Practice]{Frag mich! — the chapter, run cold}

\label{lesson:GE-R33-frag-mich}



\begin{quote}
Two sentence shapes put the verb first, and they mean different things. Name both. Then say the six W-words in one breath.
\end{quote}

\subsection*{Your turn: ask about the people and the clock}
Say these aloud:

\begin{itemize}
\raggedright
  \item \textquotedblleft{}Wer ist das?\textquotedblright{} about each — \textquotedblleft{}der Bruder, die Schwester, die Mutter, der Vater\textquotedblright{}, and the four Latin cousins they keep: \emph{frater, soror, mater, pater}
  \item \textquotedblleft{}Wann?\textquotedblright{} answered by the clock — \textquotedblleft{}Um ein Uhr\textquotedblright{}, and the Latin word inside \emph{Uhr}: \emph{hora}
\end{itemize}

\subsection*{Your turn: the distant band — flip every statement you own}
Every statement in this book can now be a question. Take the verb and move it left.

Say these aloud:

\begin{itemize}
\raggedright
  \item \textquotedblleft{}Du wohnst in Berlin.\textquotedblright{} $\to$ \textquotedblleft{}Wohnst du in Berlin?\textquotedblright{}; \textquotedblleft{}Du lernst Deutsch.\textquotedblright{} $\to$ \textquotedblleft{}Lernst du Deutsch?\textquotedblright{}
  \item the six things to eat and drink, asked for — \textquotedblleft{}Hast du Kaffee, Tee, Milch, Brot, Wasser, Wein?\textquotedblright{}
  \item the two animals, then last chapter's answer — \textquotedblleft{}Hast du einen Hund?\textquotedblright{} \textquotedblleft{}Nein, ich habe keinen Hund.\textquotedblright{} \textquotedblleft{}Hast du eine Katze?\textquotedblright{} \textquotedblleft{}Nein, ich habe keine Katze.\textquotedblright{}
  \item which of those words German bought and which it always had — \emph{Kaffee} and \emph{Tee} came from far away, \emph{Milch} and \emph{Brot} did not
\end{itemize}

\subsection*{Your turn: commands}
Strip the \emph{-en}, then say it.

Say these aloud:

\begin{itemize}
\raggedright
  \item \textquotedblleft{}Komm! Geh! Sag!\textquotedblright{}
  \item \textquotedblleft{}Hör! Öffne! Schließ!\textquotedblright{}
  \item \textquotedblleft{}Schreib!\textquotedblright{}
  \item the two whose vowel breaks — \textquotedblleft{}Lies! Sieh!\textquotedblright{}
\end{itemize}

\begin{tcolorbox}[breakable,colback=teal!4,colframe=teal!35!black,title={Before you move on}]
What turns a statement into a yes-or-no question? (\textbf{The verb moves to first place}.) What does a command leave out? (\textbf{The pronoun}.) Which sound law lines German \emph{w-} up with English \emph{wh-}? (\textbf{Grimm's}.) And which of the six W-words is two words welded together? (\emph{\textbf{Warum}} — an old \emph{was} plus \emph{um}.)
\end{tcolorbox}
